\Article{25}
To illustrate the other part of our theory as well by an example, we investigate the precision with which the side Falkenberg–Breithorn is determined from the side Wilsede–Wulfsode with the aid of adjusted observations. The function \(u\) by which the former is expressed in this case is
\SourceLeaf{097}
\[
u=22877^{\mathrm m}94\times\frac{\sin(v^{(13)}-v^{(12)}-0''652).\sin(v^{(14)}-v^{(16)}-0''814)}{\sin(v^{(1)}-v^{(0)}-0''652).\sin(v^{(6)}-v^{(4)}-0''814)}
\]

Its value, obtained from the corrected values of the directions \(v^{(0)},v^{(1)}\) etc., is
\[
=26766^{\mathrm m}68
\]

Differentiating that expression gives the following, if the differentials \(\mathrm dv^{(0)},\allowbreak\mathrm dv^{(1)}\) etc. are taken to be expressed in seconds of arc:
\[
\begin{aligned}
\mathrm du={}&0^{\mathrm m}16991(\mathrm dv^{(0)}-\mathrm dv^{(1)})+0^{\mathrm m}08836(\mathrm dv^{(4)}-\mathrm dv^{(6)})\\
&-0^{\mathrm m}03899(\mathrm dv^{(12)}-\mathrm dv^{(13)})+0^{\mathrm m}16731(\mathrm dv^{(14)}-\mathrm dv^{(16)})
\end{aligned}
\]

From this we further find
\[
\begin{aligned}\relax
[al]&=-0{,}08836\\
[bl]&=+0{,}13092\\
[cl]&=-0{,}00260\\
[dl]&=+0{,}07895\\
[el]&=+0{,}03899\\
[fl]&=-40{,}1315\\
[gl]&=+10{,}9957\\
[ll]&=+0{,}13238
\end{aligned}
\]

Finally, by the methods given above, taking the metre as the unit of linear dimensions, we find
\[
\frac1P=0{,}08329,\quad\text{or}\quad P=12{,}006
\]
whence the mean error to be expected in the value of the side Falkenberg–Breithorn is \(=0{,}2886m\) metres (where \(m\) is the mean error to be expected in the observed directions, expressed in seconds of arc); accordingly, if we adopt the value of \(m\) derived above,
\[
=0^{\mathrm m}1209
\]

Moreover, inspection of the system of triangles readily shows that the point Hauselberg could have been omitted from it altogether, while the connection between the sides Wilsede–Wulfsode and Falkenberg–Breithorn remained intact. But it is contrary to sound method to \textit{discard} observations on that account when they relate to the point Hauselberg\SourceLeaf{098}\footnote{Most of these observations had already been made before the point Breithorn was discovered and incorporated into the system.}, since they can certainly contribute to increased precision. To make clearer how much increase in precision results from this, we repeated the calculation, excluding everything that relates to the point Hauselberg; in this way eight of the 18 directions given above are removed, and the most probable errors of the remainder are found as follows:
\[
\begin{array}{r@{\;}l@{\qquad}|@{\qquad}r@{\;}l}
(0)=&+0''327&(12)=&+0''206\\
(1)=&-0{,}206&(13)=&-0{,}206\\
(3)=&-0{,}121&(14)=&+0{,}327\\
(4)=&+0{,}121&(15)=&+0{,}206\\
(6)=&-0{,}121&(16)=&+0{,}121
\end{array}
\]
The value of the side Falkenberg–Breithorn then comes out to be \(=26766^{\mathrm m}63\), differing little from the value derived above; the calculation of the weight, however, gives
\[
\frac1P=0{,}13082\quad\text{or}\quad P=7{,}644
\]
and hence the mean error to be expected is \(=0{,}36169m\) metres \(=0^{\mathrm m}1515\). It is therefore evident that, through the addition of observations; which relate to the point Hauselberg, the weight of the determination of the side Falkenberg–Breithorn has increased in the ratio of 7,644 to 12,006, or of unity to 1.571.
\par\medskip\begin{center}\rule{.28\textwidth}{.3pt}\end{center}
